% !TeX root = ../main.tex

\chapter{Dynamics Model}\label{chapter:dynamics-model}

This chapter formulates the six-degree-of-freedom dynamics used throughout the thesis. The model serves two purposes: (i) to quantify the aerodynamic forces and moments that act as disturbances on the multirotor, and (ii) to establish the reference against which the INDI controller must reject these disturbances. As discussed in Chapter~\ref{chapter:related-work}, INDI-based controllers have demonstrated strong robustness to aerodynamic loading in both classical quadrotors and tailsitter pre-transition regimes, enabling the present work to use a simplified aerodynamic description without compromising control fidelity. The aerodynamic terms introduced here are used only for offline analysis; the controller in Chapter~\ref{chapter:control-architecture} deliberately ignores them.

\section{List of Symbols}
\begin{longtable}{p{0.15\textwidth}p{0.6\textwidth}p{0.15\textwidth}}
\caption{Symbols used in the dynamics model.}
\label{tab:symbols-dynamics}\\
\toprule
\textbf{Symbol} & \textbf{Description} & \textbf{Unit} \\
\midrule
\endfirsthead
\multicolumn{3}{c}{\tablename\ \thetable\ -- continued from previous page} \\
\toprule
\textbf{Symbol} & \textbf{Description} & \textbf{Unit} \\
\midrule
\endhead
\midrule
\multicolumn{3}{r}{Continued on next page} \\
\endfoot
\bottomrule
\endlastfoot
$\mathcal{W}$ & World (inertial) frame (ENU) & --- \\
$\mathcal{B}$ & Body-fixed frame at CoG & --- \\
$\mathbf{e}_x,\mathbf{e}_y,\mathbf{e}_z$ & Body-frame unit vectors & --- \\
$\mathbf{p}$ & Position in world frame & m \\
$\mathbf{v}$ & Velocity in world frame & m/s \\
$R$ & Rotation matrix $\mathcal{B}\!\to\!\mathcal{W}$ & --- \\
$\boldsymbol{\omega}$ & Angular velocity in body frame & rad/s \\
$\mathbf{g}$ & Gravity vector $[0,0,-g]^\top$ & m/s$^2$ \\
$m$ & Vehicle mass & kg \\
$J$ & Inertia matrix in body frame & kg·m$^2$ \\
$\mathbf{F}_g$ & Gravitational force & N \\
\midrule
$f_i$ & Thrust produced by rotor $i$ & N \\
$\mathbf{u}$ & Control input vector $[f_1,f_2,f_3,f_4]^\top$ & N \\
$\omega_i$ & Angular speed of rotor $i$ & rad/s \\
$k_t$ & Thrust coefficient & N·s$^2$·rad$^{-2}$ \\
$b_i$ & Static bias in thrust model for rotor $i$ & N \\
$k_Q$ & Torque coefficient & N·m·s$^2$·rad$^{-2}$ \\
$Q_i$ & Reaction torque of rotor $i$ & N·m \\
$c_i$ & Prop rotation sign ($\pm 1$) & --- \\
\midrule
$\mathbf{v}_w$ & Wind velocity in world frame & m/s \\
$\mathbf{v}_a$ & Airspeed in body frame & m/s \\
$V$ & Airspeed magnitude & m/s \\
$\alpha$ & Angle of attack & rad \\
$\beta$ & Sideslip angle & rad \\
$\rho$ & Air density & kg/m$^3$ \\
$q$ & Dynamic pressure $=\tfrac{1}{2}\rho V^2$ & N/m$^2$ \\
$S$ & Wing reference area & m$^2$ \\
$b$ & Wing span & m \\
$AR$ & Aspect ratio $b^2/S$ & --- \\
$a_\alpha$ & Lift-curve slope & rad$^{-1}$ \\
$a_\beta$ & Side-force slope (or $C_{Y_\beta}$) & rad$^{-1}$ \\
$C_L, C_D, C_Y$ & Lift, drag, side-force coefficients & --- \\
$C_{D0}$ & Zero-lift drag coefficient & --- \\
$k$ & Induced drag factor & --- \\
\midrule
$\mathbf{F}$ & Total force in world frame & N \\
$\mathbf{F}_T$ & Rotor thrust force in world frame & N \\
$\mathbf{F}_{\!aero}$ & Aerodynamic force in body frame & N \\
$\boldsymbol{\tau}$ & Total moment in body frame & N·m \\
$\boldsymbol{\tau}_T$ & Rotor-generated moment & N·m \\
$\boldsymbol{\tau}_{\!aero}$ & Aerodynamic moment & N·m \\
$\mathbf{r}_i$ & Arm vector to rotor $i$ & m \\
$\mathbf{r}_w$ & Aerodynamic center offset & m \\
$\hat{\mathbf{t}}_i$ & Rotor $i$ thrust direction unit vector & --- \\
$G$ & Allocation matrix (force–moment mapping) & --- \\
$\hat{\mathbf{v}}$ & Airspeed unit vector & --- \\
$\hat{\mathbf{z}}_L$ & Lift direction unit vector & --- \\
$\hat{\mathbf{y}}_S$ & Side-force direction unit vector & --- \\
\end{longtable}

\section{Frames, States, Inputs}
We use a world (inertial) frame $\mathcal{W}$ with ENU convention and a body-fixed frame $\mathcal{B}$ at the CoG. Gravity is $\mathbf{g}=[0,\,0,\,-g]^\top$ in $\mathcal{W}$.

State:
\[
(\mathbf{p},\mathbf{v},R,\boldsymbol{\omega})
\in \mathbb{R}^3\times\mathbb{R}^3\times SO(3)\times\mathbb{R}^3,
\]
where $\mathbf{p},\mathbf{v}$ are position/velocity in $\mathcal{W}$, $R$ maps $\mathcal{B}\!\to\!\mathcal{W}$, and $\boldsymbol{\omega}$ is the body angular rate in $\mathcal{B}$.

Control input (per rotor $i\in\{1,\dots,4\}$):
\[
\mathbf{u}=[f_1,\,f_2,\,f_3,\,f_4]^\top, \qquad f_i\ge 0,
\]
with nominal thrust along $+\mathbf{e}_3$ of $\mathcal{B}$. Motor dynamics and allocation are covered in Chapter~\ref{chapter:control-architecture}.

\paragraph{Wind and airspeed.}
Let $\mathbf{v}_w\in\mathbb{R}^3$ be the wind in $\mathcal{W}$. Body-frame airspeed:
\[
\mathbf{v}_a \;=\; R^\top(\mathbf{v}-\mathbf{v}_w) \;\in\; \mathbb{R}^3, \quad
V \;=\;\|\mathbf{v}_a\|.
\]
Angle of attack and sideslip:
\[
\alpha = \arctan2(-v_{a,z}, v_{a,x}), \qquad
\beta = \arcsin\left(\frac{v_{a,y}}{V}\right).
\]

\section{Forces in the World Frame}
Total force:
\begin{equation}
\mathbf{F} \;=\; m\mathbf{g} \;+\; \mathbf{F}_T \;+\; R\,\mathbf{F}_{\!aero}.
\end{equation}

\subsection{Gravity}
$\mathbf{F}_g=m\mathbf{g}$.

\subsection{Rotor Thrust}
Each rotor $i$ produces a thrust $f_i$ along its (possibly tilted) direction $\hat{\mathbf{t}}_i$ in $\mathcal{B}$. The net thrust in $\mathcal{W}$ is
\begin{equation}
\mathbf{F}_T = R\sum_{i=1}^4 f_i\,\hat{\mathbf{t}}_i.
\end{equation}
For small fixed tilts (e.g., $\pm 5^\circ$) we also provide the common approximation $\mathbf{F}_T \approx (\sum_i f_i)\,R\mathbf{e}_3$; all tilt effects are embedded in the allocation matrix $G$ used for moments.

\subsection{Aerodynamic Forces on Fixed Surfaces}
We assume the wing span aligns with $\mathbf{e}_y$, so lift lies in the $(x,z)$ plane.
The definitions of $\hat{\mathbf{z}}_L$ and $\hat{\mathbf{y}}_S$ are undefined when $\hat{\mathbf{v}}$ is parallel to the respective cross-product axis (e.g., $\hat{\mathbf{v}}\parallel\mathbf{e}_y$ or $\hat{\mathbf{v}}\parallel\mathbf{e}_z$); in implementation we fall back to the last valid direction or a small-angle approximation.

Dynamic pressure $q=\tfrac{1}{2}\rho V^2$, total reference area $S$.

Unit aerodynamic directions in $\mathcal{B}$ (from $\mathbf{v}_a$):
\[
\hat{\mathbf{v}}=\frac{\mathbf{v}_a}{V},\quad
\hat{\mathbf{z}}_L=\frac{\hat{\mathbf{v}}\times\mathbf{e}_y}{\|\hat{\mathbf{v}}\times\mathbf{e}_y\|},\quad
\hat{\mathbf{y}}_S=\frac{\mathbf{e}_z\times\hat{\mathbf{v}}}{\|\mathbf{e}_z\times\hat{\mathbf{v}}\|},
\]
so that lift is perpendicular to the flow in the body $x$–$z$ plane, drag opposes $\hat{\mathbf{v}}$, and side force is lateral.

Aerodynamic force in $\mathcal{B}$:
\begin{equation}
\mathbf{F}_{\!aero}
= \underbrace{q S\,C_L(\alpha)\,\hat{\mathbf{z}}_L}_{\text{lift}}
\;-\; \underbrace{q S\,C_D(\alpha)\,\hat{\mathbf{v}}}_{\text{drag}}
\;+\; \underbrace{q S\,C_Y(\beta)\,\hat{\mathbf{y}}_S}_{\text{side}}.
\label{eq:Faero}
\end{equation}
The aerodynamic forces in~\eqref{eq:Faero} are used only for off-line analysis; they are not used in the controller.

\paragraph{Controller relevance.}
These aerodynamic forces are \emph{not} modeled in the controller; they are treated as disturbances. INDI cancels their quasi-steady contribution via incremental updates (Chapter~\ref{chapter:control-architecture}).

\section{Moments in the Body Frame}
Total moment:
\begin{equation}
\boldsymbol{\tau} \;=\; \boldsymbol{\tau}_T \;+\; \boldsymbol{\tau}_{\!aero}.
\end{equation}

\subsection{Rotor-Generated Moments}
Let $\mathbf{r}_i$ be the arm vector from CoG to rotor $i$, and $c_i\in\{+1,-1\}$ the prop rotation sign (reaction torque direction). Here $c_i\in\{+1,-1\}$ encodes the rotor spin direction (positive for CCW viewed from above). With small motor tilt handled by per-rotor thrust directions $\hat{\mathbf{t}}_i$,
\begin{align}
\boldsymbol{\tau}_{\text{arms}} &= \sum_{i=1}^4 \mathbf{r}_i \times (f_i\,\hat{\mathbf{t}}_i),\\
\boldsymbol{\tau}_{\text{react}} &= \sum_{i=1}^4 c_i\,Q_i\,\mathbf{e}_3,\qquad Q_i\approx k_Q\,\omega_i^2,
\end{align}
so that $\boldsymbol{\tau}_T=\boldsymbol{\tau}_{\text{arms}}+\boldsymbol{\tau}_{\text{react}}$. In nominal modeling $\hat{\mathbf{t}}_i\!=\!\mathbf{e}_3$; with $\pm5^\circ$ motor tilt we precompute $\hat{\mathbf{t}}_i$ and absorb it in the allocation matrix $G$ (see Chapter~\ref{chapter:control-architecture}).

\subsection{Aerodynamic Moments}
Wing forces act at an effective aerodynamic center offset $\mathbf{r}_w$ (in $\mathcal{B}$). We lump unsteady effects into a disturbance:
\begin{equation}
\boldsymbol{\tau}_{\!aero} \;=\; \mathbf{r}_w \times \mathbf{F}_{\!aero} \;+\; \boldsymbol{\tau}_{\!aero,dist}.
\end{equation}
As with forces, the controller does not model $\boldsymbol{\tau}_{\!aero}$ explicitly; INDI compensates the quasi-steady part incrementally.

\section{Equations of Motion}
Rigid-body dynamics:
\begin{align}
\dot{\mathbf{p}} &= \mathbf{v},\\
\dot{\mathbf{v}} &= \tfrac{1}{m}\big(m\mathbf{g} + \mathbf{F}_T + R\,\mathbf{F}_{\!aero}\big),\\
\dot{R} &= R[\boldsymbol{\omega}]_\times,\\
J\dot{\boldsymbol{\omega}} &= -\boldsymbol{\omega}\times J\boldsymbol{\omega} + \boldsymbol{\tau},
\end{align}
with inertia $J$ in $\mathcal{B}$ and $[\cdot]_\times$ the skew operator.

\section{Thrust and Allocation (Static Mapping)}
For simulation and identification we use static rotor maps
\[
f_i \;=\; k_t\,\omega_i^2 + b_i,\qquad Q_i \;=\; k_Q\,\omega_i^2,
\]
and a (geometry-dependent) allocation matrix $G\in\mathbb{R}^{4\times 4}$ such that
\[
\begin{bmatrix} f_T\\ \boldsymbol{\tau} \end{bmatrix}
\;=\; G\,\mathbf{f},\qquad
\mathbf{f}=[f_1,f_2,f_3,f_4]^\top,\quad f_T=\mathbf{1}^\top\mathbf{f}.
\]
Motor tilt is embedded in $G$ via $\hat{\mathbf{t}}_i$ and the arm geometry. Chapter~\ref{chapter:control-architecture} details the inverse mapping and constraints.

\subsection{Rotor Power Scaling}\label{sec:rotor-power}
Understanding how rotor power depends on rotational speed is essential for evaluating the efficiency of hybrid multirotor–fixed-wing configurations. In hover, this relationship can be derived from momentum theory, which models the rotor as an ideal actuator disk accelerating air downward to produce thrust.

For a rotor in hover, the thrust $T$ is generated by imparting a downward momentum to the air mass flow $\dot{m}$:
\[
T = \dot{m} v_i = (\rho A v_i) v_i = \rho A v_i^2,
\]
where $\rho$ is the air density, $A$ is the rotor disk area, and $v_i$ is the induced velocity through the disk. The corresponding mechanical power required to sustain this thrust is
\[
P = T v_i = \rho A v_i^3.
\]
Eliminating $v_i$ using the first equation yields the classical induced power expression from actuator-disk theory:
\[
P = T \sqrt{\frac{T}{2\rho A}} = \frac{T^{3/2}}{\sqrt{2\rho A}}
\]
(see~\cite{faa_hfh_2023}).

If blade geometry and collective pitch remain constant, blade-element theory shows that the thrust coefficient $C_T$ remains approximately constant, implying
\[
T \propto \Omega^2,
\]
where $\Omega$ is the rotor angular velocity. Substituting this into the induced power relation gives the cubic power law
\[
P \propto \Omega^3.
\]
Thus, reducing rotor speed by 20\% results in a power reduction of roughly $(0.8)^3 \approx 0.51$, or about 49\% lower power consumption—an important efficiency benefit when part of the lift is provided by wings.

This relationship strictly applies to hover or near-hover conditions, assuming uniform inflow and negligible profile and parasitic losses. In forward or edgewise flight, blade-element momentum theory (BEMT) or CFD analyses are required for accurate power prediction, but the cubic scaling remains a useful first-order engineering approximation for small deviations around hover.

The cubic law provides a theoretical baseline later compared against measured electrical-power data from outdoor experiments (Chapter~7), enabling direct verification of aerodynamic efficiency gains.

\section{Assumptions and Simplifications}
\begin{itemize}
  \item Quasi-steady aerodynamics.
  \item No explicit prop-wash/wing interference; its net effect is absorbed in the disturbance terms and identified coefficients.
  \item Rigid airframe, symmetric mass properties about $\mathcal{B}$ axes.
  \item Fixed-pitch rotors; flapping and gyroscopic coupling are neglected in the baseline model.
\end{itemize}

\section{Summary and Outlook}
The presented six-degree-of-freedom equations define a consistent reference for the vehicle's translational and rotational dynamics. Aerodynamic terms were included only to quantify expected disturbance magnitudes---order of lift, drag, and pitching moments---so that the later control analysis can assess whether these effects are sufficiently slow and quasi-steady to be rejected by INDI. No aerodynamic coefficients or force models are used online; they serve purely for interpretation and comparison.

With the governing relationships and notation established, the next step is to translate this abstract model into a physical platform whose geometry and mass distribution match the assumed parameters. The following chapter therefore details the mechanical design and manufacturing of the X-wing quadrotor, showing how the chosen dimensions, airfoil, and dihedral angles realize the modeled reference frame in practice.
